\documentclass[11pt]{article}
\usepackage{amsmath,amssymb,amsthm,mathtools,enumitem,booktabs}
\usepackage[margin=1in]{geometry}

\newtheorem{theorem}{Theorem}
\newtheorem{definition}{Definition}

\newcommand{\XiBC}{\Xi_{\mathrm{BC}}}
\newcommand{\XiHecke}{\Xi_{\mathrm{BC,Hecke}}}
\newcommand{\Psym}{\mathsf P_\infty}
\newcommand{\Pinf}{P_\infty}
\newcommand{\Mz}{M_\zeta}
\newcommand{\Ml}{M_{m_\ell}}
\newcommand{\Cg}{\mathcal Cg}
\newcommand{\Blegal}{\mathcal B}

\title{Step 172: Shifted Co-Poisson Cascade Reduction Theorem}
\author{Six Birds Foundations III}
\date{2026-05-15}

\begin{document}
\maketitle

\section{Inherited Carrier and Active Residual}

The active residual remains
\[
  \XiBC=B^*\Pi_YB
\]
on the inherited Burnol/Sonine pulled-evaluator carrier.  The sibling Hecke
residual \(\XiHecke\) remains a separate parent residual from Step 167, and
the Step 168 H6 zeta-fiber descent boundary remains \(V\)-NC
non-comparable under inherited records.

The objects used below are inherited: the carrier \(H\), the Sonin/prolate
projection \(\Pinf\), the transported Mellin projection \(\Psym\), the
Burnol generator class \(\Blegal\), the log-shift multiplier \(m_\ell\), the
transport \(T_a\), and the zero-evaluator atoms \(y_{\rho,k}\).

\section{Cascade Reduction Theorem}

\begin{theorem}[Shifted Co-Poisson Cascade Reduction for \(\XiBC\) on the Burnol/Sonine Carrier]
Under the inherited operator definitions, the Branch C shifted co-Poisson
exact route reduces as follows.

\begin{enumerate}[label=(\alph*),leftmargin=*]
\item \textbf{Step 169.} On the subclass \(u=\Cg\), with \(g\) legal and
\(G=\widehat g\) from Step 119,
\[
  M(B_{\ell,a}\Cg)
  =
  T_a\Psym\Mz(m_\ell G)
  -
  T_a\Psym\Ml\Psym\Mz G.
\]

\item \textbf{Steps 169 and 170.} The identity
\[
  M(B_{\ell,a}\Cg)=\zeta\cdot\alpha_{\ell,a,\Cg}
\]
on the full \(\Cg\) subclass is obtained if and only if the inherited
\(\mathrm{ZI\mbox{-}COV}_{\ell,a}^{\mathrm{CP}}\) obligations hold:
\[
  \Psym\Mz G=\Mz A_\infty G,
  \qquad
  T_a\Psym\Mz H=\Mz A_{T,a}H
\]
for bounded \(A_\infty\) and \(A_{T,a}\) on the typed profiles.  Under these
identities,
\[
  \alpha_{\ell,a,\Cg}
  =
  A_{T,a}\bigl(m_\ell(G-A_\infty G)\bigr).
\]

\item \textbf{Step 170.} \(\mathrm{ZI\mbox{-}COV}(i)\) holds on the
zero-free-output subclass with
\[
  A_\infty^\Omega=M_{1/\zeta}\Psym\Mz
\]
where \(1/\zeta\) is bounded on the output support.  The full-carrier
\(\mathrm{ZI\mbox{-}COV}(i)\) statement is not target-equivalent to RH under
the inherited Sonin/prolate records.

\item \textbf{Step 171.} The full-carrier gate reduces to deciding, for each
zeta-zero label \((\rho,k)\),
\[
  L_{\rho,k}(G)
  :=
  \partial_s^k(\Psym\Mz G)(\rho)
  =
  \langle \Psym\Mz G,y_{\rho,k}\rangle
  =
  \langle \Mz G,\Psym y_{\rho,k}\rangle .
\]
The five structural routes U1a--d and U2 all reduce to this same
projected-evaluator matrix element.

\item \textbf{Step 172 compositional conclusion.} Therefore closing \(\XiBC\) via the
shifted co-Poisson exact route reduces, under inherited records, to deciding
the family \(\{L_{\rho,k}\}\), or equivalently supplying full-carrier
\(\mathrm{ZI\mbox{-}COV}_{\ell,a}^{\mathrm{CP}}\), direct jet-surjectivity, or
the Mellin-side Sonin/prolate range-evaluator matrix.
\end{enumerate}
\end{theorem}

\begin{proof}
Clause (a) is the Step 169 substitution of the Step 119 co-Poisson Mellin
identity into the shifted packet formula.  Clause (b) is the Step 169 named
\(\mathrm{ZI\mbox{-}COV}_{\ell,a}^{\mathrm{CP}}\) reduction, with the Step 170
operator interpretation of \(\mathrm{ZI\mbox{-}COV}(i)\).  Clause (c) is the
Step 170 zero-free-output subclass theorem and its target-equivalence audit.
Clause (d) is the Step 171 projected zero-jet structural audit.  Clause (e) is
only their composition.  No new analytical identity is introduced here.
\end{proof}

\section{Typed External-Content Interface}

\subsection*{Branch C terminal matrix interface}

The terminal external object is
\[
  L_{\rho,k}(G)=\langle \Mz G,\Psym y_{\rho,k}\rangle.
\]
External work must respect the Burnol/Sonine carrier, legal Step 119
profiles \(G=\widehat g\), the inherited Sonin/prolate projection, and the
zero-evaluator context.  The obligation is per zeta-zero label:
\[
  \rho,\qquad 0\le k<m_\rho.
\]
The work must decide the pairing on legal \(G\), or supply an exact
Mellin-side kernel \(K_\infty(s,s')\) or equivalent range-evaluator theorem.

\subsection*{Full ZI-COV and direct jet lanes}

An alternative external route is a full-carrier proof of
\[
  \Psym\Mz G=\Mz A_\infty G,
  \qquad
  T_a\Psym\Mz H=\Mz A_{T,a}H
\]
with bounded operators in the typed carrier.  A direct jet lane may instead
prove projected jet-surjectivity or a structural refutation at explicit
\((\rho,k)\).  Such work must not replace the Sonin/prolate projection by a
zeta-zero spectral projection.

\subsection*{Sibling Burnol/Sonine branches}

Branch A remains the Calkin bridge obligation G2--G5 in \(A_\eta/K_\eta\):
faithful symbol, normal form \(C_\ell=U_\ell W_\ell+K_\ell\),
lower-faithfulness of \(U_\ell\), and compact remainder \(K_\ell\in K(H)\).

Branch B remains the per-zero finite-carrier obligation SL164.1: projected
Burnol/Sonine kernel evaluator records, plus the active matrix-source
decision.  Support-only evaluator data is not accepted bridge content.

\subsection*{Hecke sibling}

The Hecke cascade remains on \(\XiHecke\), not on \(\XiBC\).  Hecke H1--H5
remain Hecke-native obligations: direct-integral Schur legality, source
lower-frame plus Plancherel tail, Hecke Calkin bridge, per-character finite
diagnostics, and auxiliary explicit-formula records.  The H6 bridge boundary
is retained: scalar \(L(s,1)=\zeta(s)\) is not carrier identity.

\section{Cascade Terminus}

This is cascade-level diagnostic completion at the reduction-chain level for
Branch C after Steps 169--171.  It is distinct from Step 166, which summarized
the three Branch A/B/C sibling branches.  Step 172 packages the internal
Branch C reduction chain and its external-content interface.

\section{Foundational Typed-Condition Candidate}

\begin{definition}[Carrier-Typed Matrix-Element Terminality]
An adequacy residual closure problem on a typed carrier exhibits
Carrier-Typed Matrix-Element Terminality when named operator-algebra reductions
convert the closure question to a family of matrix-element pairings
\[
  \langle XG,Py\rangle
\]
whose decision is external to inherited records and logically independent of
the target conclusion unless a carrier-faithful range/evaluator theorem is
supplied.
\end{definition}

Steps 169--171 instantiate this pattern with
\[
  X=\Mz,\qquad P=\Psym,\qquad y=y_{\rho,k}.
\]
This is recorded as a candidate for foundational corpus consideration.

\section{Nonclaim Boundary}

This step does not prove RH.  It does not prove \(\XiBC=0\).  It does not
close Branch C.  It does not decide \(L_{\rho,k}\).  It does not supply
full-carrier \(\mathrm{ZI\mbox{-}COV}\).  It does not prove
jet-surjectivity.  It does not remove any retained no-go.  It does not close
the Hecke cascade.  It does not bridge H6 \(V\)-NC.  It does not begin any
Step 173 lane.

\end{document}
